%----------------------------------------------------------------------------------------
% DE ELLIPS. Alle opgaven zijn eigen werk.
%
% Opbouw: Probleem, Start, leerstof, en op een nieuwe pagina de Opdrachten, Koppel en Test jezelf.
% Labels van dit hoofdstuk beginnen met h3-.
%----------------------------------------------------------------------------------------

\mychapter{3}{De ellips}

\begin{probleem}
\naast{Kepler ontdekte in 1609 dat de aarde in een ellips rond de zon draait, met de zon in een brandpunt. Begin januari staat de aarde het dichtst bij de zon, op 147,1 miljoen km. Begin juli staat ze het verst, op 152,1 miljoen km. Hoe plat is die ellips? Als je de baan op schaal tekent, met een lange as van 20 cm: zie je dan dat het geen cirkel is? En waar staat de zon?}{\begin{tikzpicture}[scale=.5,font=\scriptsize]
  \draw[accent,thick] (0,0) ellipse (3.2 and 2.4);
  \fill[accent!60] (2.117,0) circle (5pt);
  \fill[accent] (-2.117,0) circle (1.5pt);
  \draw[lijngrijs,dashed] (-3.2,0) -- (3.2,0);
  \fill (3.2,0) circle (2.5pt) node[right] {januari};
  \fill (-3.2,0) circle (2.5pt) node[left] {juli};
  \node[below] at (2.117,-.2) {zon};
  \node[font=\tiny,text=grijs] at (0,-2.9) {(niet op schaal)};
\end{tikzpicture}}
\end{probleem}

\lkr{GO!: WD3\_06.08.11, subdoel 2 (de ellips als meetkundige plaats) en subdoel 3 (de standaardvergelijking). KOV Complementair: krommen en meetkundige plaatsen, met parametervergelijkingen. De afleiding van de vergelijking met twee wortels is een mooie oefening in algebraïsch rekenen: laat de leerlingen eerst zelf proberen, en bespreek dan waarom je een wortel eerst alleen zet. Het bewijs van de spiegeleigenschap met de spiegeling van Heron is facultatief.}

In hoofdstuk 1 bewezen de bollen van Dandelin dat een gesloten doorsnede van een kegel een ellips is: de punten waarvan de som van de afstanden tot twee brandpunten constant is. Een tuinier tekent ze met een touw en twee piketten, en planeten volgen ze rond de zon. In dit hoofdstuk geven we de ellips een vergelijking, meten we hoe plat ze is, en zien we waarom je in sommige zalen een fluistering aan de overkant hoort.

\begin{oefeningen}[Start]
\oef{h3-start-1}

\oef{h3-verken}
\end{oefeningen}

%========================================================================================
\section{De vergelijking van een ellips}

\subsection{Namen en een slim assenstelsel}

Kies de assen weer slim: de $x$-as door de brandpunten, en de oorsprong er midden tussen. Dan is $F_1(-c, 0)$ en $F_2(c, 0)$, en noemen we de constante som $2a$, met $a > c$. De ellips is dan symmetrisch ten opzichte van beide assen: spiegel je een punt over een as, dan wisselen hooguit de twee afstanden van plaats.

Waar snijdt de ellips de assen? Op de $x$-as vind je, net als bij de tuinier in hoofdstuk 1, de punten $(\pm a, 0)$. Op de $y$-as ligt een punt $B$ even ver van beide brandpunten, dus op afstand $a$ van elk. Noem zijn hoogte $b$. De driehoek $OF_2B$ is rechthoekig, en Pythagoras geeft het belangrijkste verband van dit hoofdstuk: $b^2 + c^2 = a^2$.

\begin{theorie}
\begin{minipage}[t]{.6\linewidth}
\vspace{0pt}
Een ellips met brandpunten $F_1(-c, 0)$ en $F_2(c, 0)$ en som $2a$ heeft
\begin{itemize}
\item als \textbf{\textcolor{accent}{middelpunt}} de oorsprong;
\item vier \textbf{\textcolor{accent}{toppen}}: $(\pm a, 0)$ en $(0, \pm b)$, met $\formule{b^2 = a^2 - c^2}$;
\item een \textbf{\textcolor{accent}{lange as}} van lengte $2a$ en een \textbf{\textcolor{accent}{korte as}} van lengte $2b$.
\end{itemize}
\end{minipage}\hfill
\begin{minipage}[t]{.36\linewidth}
\vspace{0pt}
\centering
\begin{tikzpicture}[scale=.62,font=\scriptsize]
  \draw[lijngrijs] (-3.4,0) -- (3.4,0) (0,-2.4) -- (0,2.5);
  \ellips{0}{0}{3}{2}
  \draw[grijs] (2.236,0) -- node[above right,inner sep=1pt] {$a$} (0,2);
  \node[left] at (0,1) {$b$};
  \node[below] at (1.118,0) {$c$};
  \brandpunt{-2.236}{0}{F_1}
  \brandpunt{2.236}{0}{F_2}
  \fill (3,0) circle (1.5pt);
  \fill (-3,0) circle (1.5pt);
  \fill (0,2) circle (1.5pt) node[above left,inner sep=1pt] {$B$};
  \fill (0,-2) circle (1.5pt);
\end{tikzpicture}
\end{minipage}
\end{theorie}

\subsection{Twee wortels wegwerken}

Nu de vergelijking. Neem een punt $P(x, y)$. De eigenschap $|PF_1| + |PF_2| = 2a$ geeft
\[ \sqrt{(x + c)^2 + y^2} + \sqrt{(x - c)^2 + y^2} = 2a . \]
Hier stoten we op een moeilijkheid. Kwadrateer je meteen, dan krijg je een dubbel product met twee wortels, en ben je er niet vanaf. Het idee: zet eerst één wortel alleen. Dan verdwijnt hij bij het kwadrateren, en wat overblijft, bevat nog maar één wortel. Die zet je opnieuw alleen, en je kwadrateert nog eens.

\begin{theorie}
\begin{stelling}{Vergelijking van een ellips}
De ellips met brandpunten $(\pm c, 0)$ en som $2a$ heeft als vergelijking \qquad $\formule{\dfrac{x^2}{a^2} + \dfrac{y^2}{b^2} = 1}$ \quad met $b^2 = a^2 - c^2$.

\vspace{1mm}
Liggen de brandpunten op de $y$-as, in $(0, \pm c)$, dan is het \qquad $\formule{\dfrac{x^2}{b^2} + \dfrac{y^2}{a^2} = 1}$
\end{stelling}

\begin{bewijs}
Zet de tweede wortel alleen en kwadrateer:
\bewijsregel{(x + c)^2 + y^2 = 4a^2 - 4a\sqrt{(x - c)^2 + y^2} + (x - c)^2 + y^2 .}
De termen $x^2$, $y^2$ en $c^2$ vallen weg: $4cx = 4a^2 - 4a\sqrt{(x - c)^2 + y^2}$. Deel door 4, zet de wortel alleen en kwadrateer opnieuw:
\bewijsregel{a^2\bigl((x - c)^2 + y^2\bigr) = (a^2 - cx)^2 \iff (a^2 - c^2)\,x^2 + a^2 y^2 = a^2(a^2 - c^2) .}
Met $b^2 = a^2 - c^2$ wordt dat $b^2x^2 + a^2y^2 = a^2b^2$. Deel door $a^2b^2$.

Omgekeerd: ligt $P$ op die kromme, dan is $y^2 = b^2 - \tfrac{b^2}{a^2}x^2$, en
\bewijsregel{|PF_2|^2 = (x - c)^2 + b^2 - \dfrac{b^2}{a^2}x^2 = \dfrac{c^2}{a^2}x^2 - 2cx + a^2 = \Bigl(a - \dfrac ca x\Bigr)^2 .}
Omdat $|x| \leqslant a$ en $c < a$, is $a - \tfrac ca x > 0$, dus $|PF_2| = a - \tfrac ca x$. Net zo is $|PF_1| = a + \tfrac ca x$, en de som is $2a$.
\end{bewijs}
\end{theorie}

Waarom dat omgekeerde deel? Kwadrateren kan oplossingen toevoegen: uit $u = v$ volgt $u^2 = v^2$, maar niet omgekeerd. Het tweede deel toont dat elk punt van de kromme de eigenschap echt heeft. Het levert bovendien een handige formule op: de afstand van een punt van de ellips tot een brandpunt hangt alleen van $x$ af.

Hoe weet je bij een vergelijking waar de brandpunten liggen? De grootste noemer hoort bij de lange as. Staat $a^2$ onder $x^2$, dan ligt de lange as horizontaal en is de ellips liggend; staat de grootste noemer onder $y^2$, dan staat ze rechtop.

\voorbeeld{}{Bepaal de toppen en de brandpunten van de ellips $9x^2 + 25y^2 = 225$. Teken de ellips.\label{vb:h3-std}}{%
\uitleg{Eerst brengen we de vergelijking in de standaardvorm, met 1 in het rechterlid. Dan lezen we $a$ en $b$ af, en volgt $c$ uit $c^2 = a^2 - b^2$.}
\naast{\begin{stappenvrij}{5cm}{4.3cm}
\dfrac{x^2}{25} + \dfrac{y^2}{9} = 1 & Deel door 225. \\[2.5ex]
a = 5, \quad b = 3 & De grootste noemer staat onder $x^2$: de lange as ligt op de $x$-as. \\[1ex]
c = \sqrt{25 - 9} = 4 & $c^2 = a^2 - b^2$. \\[1ex]
(\pm5, 0),\ (0, \pm3),\ F(\pm4, 0) & Toppen en brandpunten.
\end{stappenvrij}}{\begin{tikzpicture}[scale=.42,font=\scriptsize]
  \rooster{-5}{1}{5}{-3}{1}{3}
  \ellips{0}{0}{5}{3}
  \brandpunt[above]{-4}{0}{F_1}
  \brandpunt[above]{4}{0}{F_2}
\end{tikzpicture}}}

\voorbeeld{}{Bepaal de toppen en de brandpunten van de ellips $\dfrac{x^2}{4} + \dfrac{y^2}{9} = 1$.\label{vb:h3-staand}}{%
\uitleg{Nu staat de grootste noemer onder $y^2$: de ellips staat rechtop, en de brandpunten liggen op de $y$-as.}
\naast{\begin{stappenvrij}{5.2cm}{4.6cm}
a = 3, \quad b = 2 & $a$ hoort altijd bij de lange as. \\[1ex]
c = \sqrt{9 - 4} = \sqrt5 & $c^2 = a^2 - b^2$. \\[1ex]
(\pm2, 0),\ (0, \pm3),\ F(0, \pm\sqrt5) & Brandpunten op de lange as.
\end{stappenvrij}}{\begin{tikzpicture}[scale=.4,font=\scriptsize]
  \rooster{-2}{1}{2}{-3}{1}{3}
  \ellips{0}{0}{2}{3}
  \brandpunt[right]{0}{2.236}{F_2}
  \brandpunt[right]{0}{-2.236}{F_1}
\end{tikzpicture}}}

\begin{werkwijze}
Zorg voor 1 in het rechterlid. De grootste noemer is $a^2$, de kleinste $b^2$; de lange as ligt bij de letter met de grootste noemer. Bereken $c = \sqrt{a^2 - b^2}$ en zet de brandpunten op de lange as, op afstand $c$ van het middelpunt.
\end{werkwijze}

\nujij{Bepaal de brandpunten van $x^2 + 4y^2 = 16$.}{$\tfrac{x^2}{16} + \tfrac{y^2}{4} = 1$: $c = \sqrt{12} = 2\sqrt3$, dus $F(\pm 2\sqrt3, 0)$.}

En omgekeerd? Ken je twee van de drie getallen $a$, $b$ en $c$, dan volgt de derde uit $b^2 + c^2 = a^2$. Dat is de rechthoekige driehoek uit de figuur.

\voorbeeld{}{Geef een vergelijking van de ellips met brandpunten $(\pm2, 0)$ die door het punt $(0, 3)$ gaat.\label{vb:h3-opstellen}}{%
\uitleg{Het punt $(0, 3)$ ligt op de $y$-as, dus het is een top op de korte as: $b = 3$. De brandpunten geven $c = 2$.}
\begin{stappen}
a^2 = b^2 + c^2 = 9 + 4 = 13 & Pythagoras in de driehoek $OF_2B$. \\[1ex]
\dfrac{x^2}{13} + \dfrac{y^2}{9} = 1 & De lange as ligt op de $x$-as, bij de brandpunten.
\end{stappen}}

\begin{oefeningen}
\oef{h3-std}
\end{oefeningen}

%========================================================================================
\section{Hoe plat is een ellips?}

\subsection{De excentriciteit}

Een ellips met $a = 5$ en $b = 4{,}9$ lijkt sterk op een cirkel, een ellips met $a = 5$ en $b = 1$ niet. Hoe druk je dat uit in één getal? Kijk naar de brandpunten. Bij een cirkel vallen ze samen in het middelpunt: $c = 0$. Hoe verder ze uit elkaar liggen ten opzichte van de lange as, hoe platter de ellips.

\begin{theorie}
De \textbf{\textcolor{accent}{excentriciteit}} van een ellips is \qquad $\formule{e = \dfrac ca}$ \qquad met $0 \leqslant e < 1$.
\end{theorie}

Bij $e = 0$ is het een cirkel. Hoe dichter $e$ bij 1, hoe platter. Want $\tfrac ba = \tfrac{\sqrt{a^2 - c^2}}{a} = \sqrt{1 - e^2}$: de verhouding van de assen hangt alleen van $e$ af. De figuur toont vier ellipsen met dezelfde lange as.

\figuren{\begin{tikzpicture}[scale=.36,font=\scriptsize]
  \ellips{0}{0}{3}{3}
  \fill[accent] (0,0) circle (3pt);
  \node at (0,-3.8) {$e = 0$};
\end{tikzpicture}\qquad
\begin{tikzpicture}[scale=.36,font=\scriptsize]
  \ellips{0}{0}{3}{2.598}
  \fill[accent] (-1.5,0) circle (3pt) (1.5,0) circle (3pt);
  \node at (0,-3.8) {$e = 0{,}5$};
\end{tikzpicture}\qquad
\begin{tikzpicture}[scale=.36,font=\scriptsize]
  \ellips{0}{0}{3}{1.8}
  \fill[accent] (-2.4,0) circle (3pt) (2.4,0) circle (3pt);
  \node at (0,-3.8) {$e = 0{,}8$};
\end{tikzpicture}\qquad
\begin{tikzpicture}[scale=.36,font=\scriptsize]
  \ellips{0}{0}{3}{.937}
  \fill[accent] (-2.85,0) circle (3pt) (2.85,0) circle (3pt);
  \node at (0,-3.8) {$e = 0{,}95$};
\end{tikzpicture}}

Opvallend: bij $e = 0{,}5$ zie je nauwelijks dat het geen cirkel is, want $b = 0{,}87a$. Maar de brandpunten liggen al ver uit het midden. Dat is precies het antwoord op het probleem.

\voorbeeld{}{Terug naar het probleem. De aarde staat op 147,1 miljoen km van de zon in januari, en op 152,1 miljoen km in juli. Bereken $a$, $c$, $b$ en $e$ van de baan. Teken je ze met een lange as van 20 cm, hoeveel scheelt de korte as dan van 20 cm, en hoe ver staat de zon uit het midden?\label{vb:h3-terug}}{%
\uitleg{De zon staat in een brandpunt. De kleinste en de grootste afstand vind je in de toppen op de lange as: $a - c$ en $a + c$.}
\begin{stappen}
a - c = 147{,}1, \quad a + c = 152{,}1 & Tel op en trek af. \\[1ex]
a = 149{,}6, \quad c = 2{,}5 & In miljoen km. \\[1ex]
e = \dfrac{2{,}5}{149{,}6} \approx 0{,}0167 & Een heel kleine excentriciteit. \\[2ex]
\dfrac ba = \sqrt{1 - e^2} \approx 0{,}99986 & Op 20 cm: de korte as is $2b \approx 19{,}997$ cm, een verschil van 0,03 mm. \\[1ex]
c = e \cdot a \approx 0{,}0167 \cdot 10 = 0{,}167 & Op schaal staat de zon 1,7 mm uit het midden.
\end{stappen}
De baan is dus op het oog een perfecte cirkel: het verschil in de assen is kleiner dan de dikte van een potloodlijn. Maar de zon staat zichtbaar uit het midden. De afstand tot de zon schommelt daardoor met ongeveer 3\%, te weinig om de seizoenen te verklaren: die komen van de schuine stand van de aardas.}

\nujij{Een ellips heeft $a = 10$ en $e = 0{,}6$. Bereken $c$ en $b$.}{$c = ea = 6$ en $b = \sqrt{100 - 36} = 8$.}

\begin{oefeningen}
\oef{h3-exc}
\end{oefeningen}

\subsection{Een samengedrukte cirkel}

Er is nog een manier om naar een ellips te kijken. Neem de cirkel $x^2 + y^2 = a^2$ en druk hem verticaal samen: vermenigvuldig elke $y$-coördinaat met $\tfrac ba$. Een punt $(x, y)$ van de ellips komt dan van het punt $\bigl(x, \tfrac ab y\bigr)$ van de cirkel, en dat voldoet aan $x^2 + \tfrac{a^2}{b^2}y^2 = a^2$. Deel door $a^2$: dat is de vergelijking van de ellips. Een ellips is dus een samengedrukte cirkel.

\naast{Dat geeft meteen een manier om de punten van een ellips te beschrijven met een hoek. Een punt van de cirkel is $(a\cos t, a\sin t)$. Samengedrukt wordt dat $(a\cos t, b\sin t)$. Let op: $t$ is de hoek bij het punt van de cirkel, niet bij het punt van de ellips.

Het levert ook de oppervlakte. Samendrukken met factor $\tfrac ba$ maakt elke verticale strook $\tfrac ba$ keer zo lang, en dus ook de hele oppervlakte: $\pi a^2 \cdot \tfrac ba = \pi ab$.}{\begin{tikzpicture}[scale=.52,font=\scriptsize]
  \draw[lijngrijs] (-3.4,0) -- (3.4,0) (0,-3.3) -- (0,3.4);
  \draw[lijngrijs,dashed] (0,0) circle (3);
  \ellips{0}{0}{3}{1.8}
  \draw[grijs] (0,0) -- (1.5,2.598);
  \draw[grijs,dashed] (1.5,2.598) -- (1.5,0);
  \fill (1.5,2.598) circle (2pt) node[above right] {$(a\cos t, a\sin t)$};
  \fill[accent] (1.5,1.559) circle (2pt) node[right,text=black] {$(a\cos t, b\sin t)$};
  \draw[grijs] (.7,0) arc[start angle=0,end angle=60,radius=.7];
  \node at (1,.38) {$t$};
\end{tikzpicture}}

\begin{theorie}
\begin{stelling}{Parametervergelijking en oppervlakte}
De punten van de ellips $\dfrac{x^2}{a^2} + \dfrac{y^2}{b^2} = 1$ zijn de punten \quad $\formule{\begin{cases} x = a\cos t \\ y = b\sin t \end{cases}}$ \quad met $t \in [0, 2\pi[$.

\vspace{1mm}
De oppervlakte van de ellips is \quad $\formule{A = \pi ab}$
\end{stelling}
\end{theorie}

Dat elk punt $(a\cos t, b\sin t)$ op de ellips ligt, controleer je door in te vullen: $\cos^2 t + \sin^2 t = 1$. De oppervlakte met integralen berekenen kan ook, zie de Koppel. Voor de omtrek bestaat er geen eenvoudige formule: dat is een van de redenen waarom de ellips wiskundigen eeuwenlang bezighield (opdracht~\ref{oef:h3-o13}).

\voorbeeld{}{Welk punt van de ellips $\dfrac{x^2}{16} + \dfrac{y^2}{9} = 1$ hoort bij $t = \dfrac{\pi}{3}$? Bereken ook de oppervlakte van de ellips.\label{vb:h3-param}}{%
\uitleg{De parametervergelijking geeft het punt rechtstreeks, met $a = 4$ en $b = 3$.}
\begin{stappen}
x = 4\cos\tfrac{\pi}{3} = 2, \quad y = 3\sin\tfrac{\pi}{3} = \tfrac{3\sqrt3}{2} & Het punt $\bigl(2, \tfrac{3\sqrt3}{2}\bigr)$. \\[1ex]
\tfrac{4}{16} + \tfrac{27/4}{9} = \tfrac14 + \tfrac34 = 1 & Controle in de vergelijking. \\[1ex]
A = \pi \cdot 4 \cdot 3 = 12\pi \approx 37{,}7 & Tussen de cirkels met straal 3 en 4: $9\pi < 12\pi < 16\pi$.
\end{stappen}}

\nujij{Bereken de oppervlakte van de ellips $x^2 + 4y^2 = 4$.}{$\tfrac{x^2}{4} + y^2 = 1$: $A = \pi \cdot 2 \cdot 1 = 2\pi$.}

\begin{oefeningen}
\oef{h3-param}
\end{oefeningen}

%========================================================================================
\section{Verschuiven en raken}

\subsection{Een verschoven ellips}

Ligt het middelpunt niet in de oorsprong, dan werkt het zoals bij de parabool: vervang $x$ door $x - x_0$ en $y$ door $y - y_0$. De ellips met middelpunt $(x_0, y_0)$ is
\[ \frac{(x - x_0)^2}{a^2} + \frac{(y - y_0)^2}{b^2} = 1 , \]
en de brandpunten schuiven mee. Werk je de haakjes uit, dan krijg je een vergelijking als $4x^2 + 9y^2 - 16x + 54y + 61 = 0$. Daarin herken je een ellips aan de kwadraten $x^2$ en $y^2$ met positieve, verschillende coëfficiënten. Om het middelpunt te vinden, splits je weer het kwadraat af.

\voorbeeld{}{Bepaal het middelpunt, de toppen en de brandpunten van de ellips $4x^2 + 9y^2 - 16x + 54y + 61 = 0$.\label{vb:h3-versch}}{%
\uitleg{Zet eerst de coëfficiënt van het kwadraat buiten haakjes, en splits dan pas het kwadraat af. Wat je toevoegt, tel je met die coëfficiënt mee.}
\naast{\begin{stappenvrij}{6.2cm}{3.8cm}
4(x^2 - 4x) + 9(y^2 + 6y) + 61 = 0 & Groepeer. \\[1ex]
4(x - 2)^2 - 16 + 9(y + 3)^2 - 81 + 61 = 0 & $4 \cdot 4 = 16$ en $9 \cdot 9 = 81$. \\[1ex]
4(x - 2)^2 + 9(y + 3)^2 = 36 & Constanten naar rechts. \\[1ex]
\dfrac{(x - 2)^2}{9} + \dfrac{(y + 3)^2}{4} = 1 & Deel door 36. \\[2.5ex]
M(2, -3),\ a = 3,\ b = 2,\ c = \sqrt5 & Liggend. \\[1ex]
F(2 \pm \sqrt5, -3) & Toppen $(-1, -3)$, $(5, -3)$, $(2, -1)$, $(2, -5)$.
\end{stappenvrij}}{\begin{tikzpicture}[scale=.5,font=\scriptsize]
  \rooster{-1}{1}{5}{-5}{1}{1}
  \ellips{2}{-3}{3}{2}
  \brandpunt[above]{-.236}{-3}{F_1}
  \brandpunt[above]{4.236}{-3}{F_2}
  \fill (2,-3) circle (2.5pt);
\end{tikzpicture}}}

\begin{werkwijze}
Zet de coëfficiënten van $x^2$ en $y^2$ buiten haakjes, splits het kwadraat af en breng de constanten naar rechts. Deel door het rechterlid, zodat er 1 staat. Lees het middelpunt af, en werk verder zoals bij een ellips met middelpunt in de oorsprong.
\end{werkwijze}

\nujij{Bepaal het middelpunt van $x^2 + 2y^2 + 2x - 8y + 1 = 0$.}{$(x + 1)^2 + 2(y - 2)^2 = 8$: middelpunt $(-1, 2)$.}

\begin{oefeningen}
\oef{h3-versch}
\end{oefeningen}

\subsection{De raaklijn}

Bij de parabool vonden we de raaklijn met de afgeleide, en vatten we het resultaat samen als ontdubbelen. Werkt dat ook hier? We gebruiken nu een ander idee, dat ook uitlegt wat een raaklijn eigenlijk doet: ze raakt de ellips in $P$, en de hele ellips ligt aan één kant van de raaklijn.

\begin{theorie}
\begin{stelling}{Raaklijn aan een ellips}
De raaklijn in het punt $P(x_1, y_1)$ van de ellips $\dfrac{x^2}{a^2} + \dfrac{y^2}{b^2} = 1$ is \qquad $\formule{\dfrac{x_1x}{a^2} + \dfrac{y_1y}{b^2} = 1}$
\end{stelling}

\begin{bewijs}
$P$ ligt op die rechte, want $\tfrac{x_1^2}{a^2} + \tfrac{y_1^2}{b^2} = 1$. Neem nu een willekeurig punt $Q(x, y)$ van de ellips. Een som van kwadraten is nooit negatief:
\bewijsregel{0 \leqslant \dfrac{(x - x_1)^2}{a^2} + \dfrac{(y - y_1)^2}{b^2} = \Bigl(\dfrac{x^2}{a^2} + \dfrac{y^2}{b^2}\Bigr) + \Bigl(\dfrac{x_1^2}{a^2} + \dfrac{y_1^2}{b^2}\Bigr) - 2\Bigl(\dfrac{x_1x}{a^2} + \dfrac{y_1y}{b^2}\Bigr) .}
De eerste twee haakjes zijn 1, omdat $Q$ en $P$ op de ellips liggen. Dus
\bewijsregel{\dfrac{x_1x}{a^2} + \dfrac{y_1y}{b^2} \leqslant 1 ,}
met gelijkheid alleen als $Q = P$. De hele ellips ligt dus aan één kant van de rechte, en raakt ze alleen in $P$.
\end{bewijs}
\end{theorie}

De formule is weer ontdubbelen: $x^2 \to x_1x$ en $y^2 \to y_1y$. Wie liever de afgeleide gebruikt, vindt hetzelfde (opdracht~\ref{oef:h3-o19}).

\voorbeeld{}{Bepaal de raaklijn aan de ellips $\dfrac{x^2}{25} + \dfrac{y^2}{9} = 1$ in het punt $P\bigl(3, \tfrac{12}{5}\bigr)$.\label{vb:h3-raak}}{%
\uitleg{Eerst controleren we dat $P$ op de ellips ligt, dan ontdubbelen we.}
\naast{\begin{stappenvrij}{4.9cm}{4.2cm}
\dfrac{9}{25} + \dfrac{144/25}{9} = \dfrac{9}{25} + \dfrac{16}{25} = 1 & $P$ ligt op de ellips. \\[2.5ex]
\dfrac{3x}{25} + \dfrac{12y}{45} = 1 & Ontdubbelen. \\[2.5ex]
9x + 20y = 75 & Vermenigvuldig met 75.
\end{stappenvrij}}{\begin{tikzpicture}[scale=.4,font=\scriptsize]
  \rooster{-5}{1}{6}{-3}{1}{4}
  \ellips{0}{0}{5}{3}
  \draw[grijs,thick] (-1,4.2) -- (7,.6);
  \fill (3,2.4) circle (3pt) node[below left] {$P$};
\end{tikzpicture}}}

\nujij{Bepaal de raaklijn aan $\dfrac{x^2}{8} + \dfrac{y^2}{2} = 1$ in $(2, 1)$.}{$\tfrac{2x}{8} + \tfrac{y}{2} = 1$, dus $x + 2y = 4$.}

\begin{oefeningen}
\oef{h3-raak}
\end{oefeningen}

\subsection{De fluistergalerij}

In sommige zalen met een elliptisch gewelf gebeurt iets vreemds: wie in het ene brandpunt fluistert, wordt in het andere brandpunt duidelijk gehoord, terwijl iemand ertussen niets verstaat. Waarom? Net als bij de parabool gaat het om weerkaatsing: geluid dat vanuit $F_1$ vertrekt, moet na één botsing door $F_2$ gaan.

Het idee is mooi. Elk punt $Q$ van de raaklijn, behalve $P$, ligt buiten de ellips, en daar is de som van de afstanden tot de brandpunten groter dan $2a$. Van alle punten van de raaklijn heeft $P$ dus de kleinste som. En de kortste weg van $F_1$ via een rechte naar $F_2$ maakt gelijke hoeken met die rechte, zoals een lichtstraal op een spiegel.

\begin{theorie}
\begin{minipage}[t]{.6\linewidth}
\vspace{0pt}
\begin{stelling}{Spiegeleigenschap}
De raaklijn in een punt $P$ van een ellips maakt gelijke hoeken met $PF_1$ en $PF_2$.

\vspace{1mm}
Een straal uit het ene brandpunt weerkaatst dus door het andere.
\end{stelling}
\end{minipage}\hfill
\begin{minipage}[t]{.36\linewidth}
\vspace{0pt}
\centering
\begin{tikzpicture}[scale=.55,font=\scriptsize]
  \ellips{0}{0}{3}{2}
  \draw[grijs,thick] (-.3,2.65) -- (3.9,.55);
  \draw[accent!70] (-2.236,0) -- (1.8,1.6) -- (2.236,0);
  \draw[lijngrijs,dashed] (1.8,1.6) -- (3.342,2.211) node[right,text=black] {$F_2'$};
  \draw[lijngrijs,dashed] (2.236,0) -- (3.342,2.211);
  \fill (3.342,2.211) circle (1.5pt);
  \fill (1.8,1.6) circle (2pt) node[above] {$P$};
  \brandpunt{-2.236}{0}{F_1}
  \brandpunt{2.236}{0}{F_2}
  \node[text=grijs,above] at (3.9,.55) {$t$};
\end{tikzpicture}
\end{minipage}

\begin{bewijs}
Ligt $Q$ buiten de ellips, dan snijdt $[QF_1]$ de ellips in een punt $R$, en
\bewijsregel{|QF_1| + |QF_2| = |RF_1| + |RQ| + |QF_2| \geqslant |RF_1| + |RF_2| = 2a .}
Zo heeft $P$ van alle punten van de raaklijn $t$ de kleinste som. Spiegel $F_2$ in $t$ tot $F_2'$. Voor elk punt $X$ van $t$ is $|XF_1| + |XF_2| = |XF_1| + |XF_2'|$, en dat is het kleinst als $X$ op $[F_1F_2']$ ligt. Dus ligt $P$ op $[F_1F_2']$. Dan maakt $PF_1$ met $t$ dezelfde hoek als $PF_2'$ (overstaande hoeken), en die is gelijk aan de hoek van $PF_2$ met $t$ (spiegeling).
\end{bewijs}
\end{theorie}

Die eigenschap wordt ook in de geneeskunde gebruikt. Een niersteenvergruizer heeft een spiegel in de vorm van een halve ellipsoïde: een schokgolf die in het ene brandpunt ontstaat, komt in het andere brandpunt samen, precies waar de niersteen ligt. Het weefsel ertussen krijgt maar een fractie van de energie.

%========================================================================================
\opdrachtenkop

\begin{oefeningen}
\oef{h3-o1}

\oef{h3-o2}

\oef{h3-o3}

\oef{h3-o4}

\oef{h3-o5}

\oef{h3-o6}

\oef{h3-o7}

\oef{h3-o8}

\oef{h3-o9}

\oef{h3-o10}

\oef{h3-o11}

\oef{h3-o12}

\oef{h3-o13}

\oef{h3-o14}

\oef{h3-o15}

\oef{h3-o16}

\oef{h3-o17}

\oef{h3-o18}

\oef{h3-o19}

\oef{h3-o20}

\oef{h3-o21}
\end{oefeningen}

%----------------------------------------------------------------------------------------
% KOPPEL: verbanden met andere leerstof
%----------------------------------------------------------------------------------------
\needspace{20\baselineskip}
\opdrachtsectie{Koppel}

\lkr{Integralen en Ruimtemeetkunde staan net als Kegelsneden in het 6de jaar: kijk na of ze al gezien zijn. Matrices zijn gezien in het 5de jaar.}

\begin{oefeningen}
\oef{h3-koppel-1}

\oef{h3-koppel-2}

\oef{h3-koppel-3}
\end{oefeningen}

%----------------------------------------------------------------------------------------
\opdrachtsectie{Test jezelf}

\begin{oefeningen}
\oef{h3-t1}

\oef{h3-t2}

\oef{h3-t3}

\oef{h3-t4}

\oef{h3-t5}
\end{oefeningen}

\herhaalschema{%
  \fv{oef:h3-t1}{vb:h3-std}{oef:h3-std}%
  \fv{oef:h3-t2}{vb:h3-opstellen}{oef:h3-o2}%
  \fv{oef:h3-t3}{vb:h3-terug}{oef:h3-exc}%
  \fv{oef:h3-t4}{vb:h3-versch}{oef:h3-versch}%
  \fv{oef:h3-t5}{vb:h3-raak}{oef:h3-raak}}

\ifzelfstudie
\section*{Maak je eigen samenvatting}
Sluit dit boek en schrijf in je schrift: de namen $a$, $b$ en $c$ met hun verband, de standaardvergelijking en hoe je de stand herkent, het idee van het bewijs met twee wortels, de excentriciteit, de parametervergelijking en de oppervlakte, de raaklijn door te ontdubbelen, en de spiegeleigenschap. Vergelijk daarna met de kaders.
\fi
